%! Setting Up the SVD — Unit 7, Lesson 7.0 — Exit Ticket (Answer Key)
\documentclass[10pt]{article}
\usepackage{linalg-article}
\usepackage{linalg-key}   % pulls in -boxes; do NOT also load -boxes
\newcommand{\TallMath}[1]{$\displaystyle #1\rule[-1.4em]{0pt}{3.2em}$}
\newcommand{\ee}{\mathbf{e}}

\begin{document}

\pageheader{Unit 7, Lesson 7.0}{Exit Ticket --- Answer Key}
\namedateperiod

\vspace{0.15in}

No notes. Show your reasoning.

\begin{enumerate}[leftmargin=1.8em, itemsep=16pt]
  \item \textbf{Read the ellipse.} A diagonal matrix $D=\begin{bmatrix} 4 & 0 \\ 0 & 3 \end{bmatrix}$
        turns the unit circle into an ellipse. What are its two half-widths, and what are the singular
        values $\sigma_1,\sigma_2$ (larger first)?
        \ansline{Half-widths $4$ (along the $x$-axis) and $3$ (along the $y$-axis), so $\sigma_1=4$ and $\sigma_2=3$.}
  \item \textbf{Perpendicular in, perpendicular out.} For $S=\begin{bmatrix} 0 & 3 \\ 2 & 0 \end{bmatrix}$,
        compute $S\ee_1$ and $S\ee_2$. Are the two outputs perpendicular? Give the singular values
        $\sigma_1,\sigma_2$.
        \ansline{$S\ee_1=\begin{bsmallmatrix} 0\\2 \end{bsmallmatrix}$ (length $2$), $S\ee_2=\begin{bsmallmatrix} 3\\0 \end{bsmallmatrix}$ (length $3$); dot $=0$, so yes, perpendicular. $\sigma_1=3$, $\sigma_2=2$.}
  \item \textbf{Explain.} In one sentence, what does a \emph{singular value} of a matrix tell you
        \emph{geometrically}?
        \ansline{It is a half-width (semi-axis length) of the ellipse the unit circle is stretched into --- how far the matrix stretches along one of its perpendicular singular directions.}
\end{enumerate}

\begin{teachernote}
Sort responses into ``got it / arithmetic slip / concept slip.'' Item~1 is the clean read-off from a
diagonal matrix ($\sigma_1=4,\sigma_2=3$). Item~2 checks the perpendicular-in/perpendicular-out idea and
the ordering convention on a non-diagonal (swap) matrix. Item~3 is the target sentence: a singular value
is a stretch factor / semi-axis length of the ellipse. If many stumble on item~3, open Lesson~7.1 with a
30-second ``circle $\to$ ellipse, singular values are the half-widths'' re-cap.
\end{teachernote}

\end{document}
